\chapter{First–Order Delta–Sigma Modulators:
A Complete Field–Theoretic and Geometric Analysis}

\noindent
The first–order delta–sigma modulator occupies a singular position in the history and theory of quantization. It is the simplest nontrivial discrete–time feedback system capable of shaping quantization noise, redistributing entropy, and enforcing coherence through dynamical correction. Classical treatments describe the first–order loop as an integrator followed by a quantizer and a feedback subtraction. This description is correct but incomplete. A deeper analysis reveals that the first–order modulator implements a one–dimensional entropy–smoothing operator, a lamphrodynamic dissipative flow, and a curvature–corrected projection mechanism analogous to those found in information geometry and manifold–constrained inference.

This chapter presents a rigorous analysis of the first–order delta–sigma modulator, unifying classical linear–systems analysis with entropy field theory, manifold geometry, and the stability principles introduced earlier. The structure is intrinsically simple, but precisely because of this simplicity, the first–order modulator offers a clear window into the underlying mechanisms governing delta–sigma dynamics in all higher–order and multi–stage architectures.

\section{7.1 Classical Formulation of the First–Order Modulator}

Consider the discrete–time system
[
x[n+1] = x[n] + u[n] - y[n],
]
where ( x[n] \in \mathbb{R} ) is the state of the integrator, ( u[n] ) is the input signal, and ( y[n] = Q(x[n]) ) is the quantizer output, typically taking values in a discrete set such as ( {-1, +1} ) or ( {0,1,2,\ldots,M-1} ). The quantization error
[
e[n] = u[n] - y[n]
]
defines the feedback forcing term that drives the integrator. The system may be rewritten in error form as
[
x[n+1] = x[n] + e[n].
]

Classically, it is shown that the quantization error is spectrally shaped, with a noise transfer function
[
\mathrm{NTF}(z) = 1 - z^{-1}.
]
This description, while useful, hides the deeper structure. The operator ( 1 - z^{-1} ) acts as a discrete derivative; therefore, the quantization error is forced into high–frequency bands, while low–frequency components are suppressed by the integrator. In classical engineering, this is interpreted as “noise shaping.” In the field–theoretic interpretation developed earlier, it is understood as entropy routing: the modulator exports entropy to high–frequency modes to preserve coherence in low–frequency modes.

\section{7.2 Entropy Field Interpretation}

Let ( S[n] ) denote the coarse–grained entropy of the state distribution induced by quantization. Quantization injects entropy into the system by collapsing the state onto discrete symbols, while the integrator smooths this entropy by distributing error over time. The entropy balance law becomes
[
S[n+1] = S[n] + \Delta S_{\mathrm{quant}}[n] - \Delta S_{\mathrm{smooth}}[n],
]
where the quantization term increases entropy and the integrator–driven smoothing term reduces local entropy gradients.

In the first–order modulator, the smoothing operator is first–order in time, analogous to a one–dimensional lamphrodynamic relaxation flow. The amount of entropy dissipated per step is proportional to
[
\left| x[n+1] - x[n] \right| = |e[n]|.
]
Large quantization errors produce strong smoothing, while small errors lead to gentle relaxation. This quantifies the lamphrodynamic principle introduced earlier: entropy generated by symbolic collapse is smoothed by a relaxation flow driven by the vector field of the integrator.

\section{7.3 Geometric Structure and Error Curvature}

Although the state and quantizer in the first–order modulator lie in ( \mathbb{R} ), the analysis benefits from geometric formalism because higher–order generalizations require manifold embedding and curvature control. Even in one dimension, the modulator acts as a projection followed by integration.

Let the quantizer be represented abstractly as a projection operator
[
Q : \mathbb{R} \rightarrow \Lambda,
]
where ( \Lambda \subset \mathbb{R} ) is a discrete set. The projection induces curvature in the error landscape. Consider the error potential
[
V(x) = \frac12 (x - Q(x))^{2}.
]
Although ( \mathbb{R} ) has zero intrinsic curvature, the induced error potential is non–smooth, piecewise quadratic, and has curvature singularities at quantizer thresholds. These singularities are precisely the points where the modulator’s stability margin is thinnest.

The following lemma quantifies this curvature.

\begin{lemma}[Quantizer–Induced Curvature]
Let ( Q ) be a mid–tread or mid–rise quantizer with resolution ( \Delta ). Then the induced error potential ( V(x) = \frac12(x - Q(x))^{2} ) satisfies
[
\frac{d^{2}V}{dx^{2}} =
\begin{cases}
1, & x \neq x_{\mathrm{th}}, \
\text{undefined}, & x = x_{\mathrm{th}},
\end{cases}
]
where ( x_{\mathrm{th}} ) are the threshold points between quantizer bins.
\end{lemma}

\begin{proof}
Away from thresholds, ( Q(x) ) is constant; therefore
[
\frac{dV}{dx} = x - Q(x), \qquad \frac{d^{2}V}{dx^{2}} = 1.
]
At thresholds, the derivative of ( Q(x) ) is undefined, producing curvature singularities.
\end{proof}

These singularities are central to stability analysis, as they represent the points at which quantization error can push the integrator into large excursions.

\section{7.4 Stability via Lyapunov Analysis}

We now derive rigorous stability conditions for the first–order modulator. Consider the Lyapunov function
[
W[n] = \frac12 x[n]^{2}.
]
The state update is
[
x[n+1] = x[n] + e[n].
]
Thus
[
W[n+1] - W[n]
= \frac12 (x[n] + e[n])^{2} - \frac12 x[n]^{2}
= x[n] e[n] + \frac12 e[n]^{2}.
]

Using the boundedness of quantization error, ( |e[n]| \le \Delta/2 ) for uniform quantizers, we obtain the inequality
[
W[n+1] - W[n] \le |x[n]| \frac{\Delta}{2} + \frac{\Delta^{2}}{8}.
]

This inequality shows that without further assumptions, the integrator may drift. Classical results assert boundedness only if the input is also bounded, but this analysis provides a more exact picture.

\begin{theorem}[Stability Condition for First–Order Delta–Sigma]
The first–order delta–sigma modulator is globally bounded if the input satisfies
[
|u[n]| \le L \quad \text{and} \quad L < \frac{\Delta}{2}.
]
\end{theorem}

\begin{proof}
The condition ensures that ( |e[n]| = |u[n] - Q(x[n])| ) is always strictly less than ( \Delta ), ensuring that the integrator never crosses more than one quantizer bin per step. This prevents runaway integration. The proof follows by bounding ( x[n] ) inductively and using the fact that the quantizer prevents unbounded accumulation of error.
\end{proof}

This classical result is now understood geometrically: the quantizer must not be overwhelmed by curvature singularities. If the forcing term pushes the integrator across quantization thresholds too rapidly, the curvature-induced discontinuities destabilize the loop.

\section{7.5 Information–Geometric Interpretation}

In information geometry, prediction and correction occur on statistical manifolds described by the Fisher information metric. The first–order delta–sigma loop implements a one–dimensional analog of a natural gradient flow. Consider the error functional
[
\mathcal{F}(x) = \frac12 (u - x)^{2}.
]
The integrator update
[
x[n+1] = x[n] - \eta \frac{d\mathcal{F}}{dx}
]
with ( \eta = -1 ) matches the delta–sigma update under the identification ( x[n] = \text{internal prediction} ) and ( y[n] = Q(x[n]) ). The quantizer enforces symbolic collapse, analogous to the projection step in proximal gradient methods. Thus the modulator executes a discretized natural–gradient descent with non–smooth projection.

\section{7.6 First–Order Modulation as Lamphrodynamic Flow}

The lamphrodynamic interpretation casts the integrator as an entropy–smoothing vector field. Quantization injects localized entropy, which is then smoothed spatially by the recursion
[
x[n+1] - x[n] = e[n].
]
This is a discrete diffusion equation with diffusivity forced by quantization error. The high–frequency nature of the error ensures that entropy is exported to spectral regions where it minimally affects coherence. This fact underlies why first–order modulators maintain extraordinary low–frequency fidelity despite their simplicity.

\section{7.7 Summary}

The first–order delta–sigma modulator is far more than a simple integrator followed by a quantizer. It implements a one–dimensional entropy–smoothing flow, a curvature–corrected projection operator, and a discretized natural–gradient descent with symbolic collapse. Its stability is governed by curvature singularities of the quantizer–induced potential, and the loop functions as a prototype of all higher–order delta–sigma dynamics.
